\section{Biomechanical Center of Gravity Modeling}
\label{sec:biomechanics}

To quantify athletic stability and stance depth from 2D coordinate streams, we implement a linked-segment human body model parameterized by de Leva's adjustments to Zatsiorsky-Seluyanov's segment inertia parameters \cite{deleva1996adjustments}.

\begin{table}[h]
\centering
\caption{Segment Mass Proportions (de Leva, 1996)}
\label{tab:deleva_weights}
\begin{tabular}{lcc}
\hline
\textbf{Segment ($i$)} & \textbf{Relative Mass ($w_i$)} & \textbf{Proximal Boundary $\to$ Distal} \\
\hline
Head / Neck & $0.0681$ & Top of Head $\to$ C7 / Acromion Midpoint \\
Trunk / Torso & $0.4346$ & Acromion Midpoint $\to$ Greater Trochanter Midpoint \\
Upper Arm ($\times 2$) & $0.0263$ & Acromion $\to$ Lateral Epicondyle \\
Forearm ($\times 2$) & $0.0150$ & Lateral Epicondyle $\to$ Ulnar Styloid \\
Hand ($\times 2$) & $0.0059$ & Ulnar Styloid $\to$ Metacarpophalangeal \\
Thigh ($\times 2$) & $0.1416$ & Greater Trochanter $\to$ Femoral Condyle \\
Shank ($\times 2$) & $0.0433$ & Femoral Condyle $\to$ Lateral Malleolus \\
Foot ($\times 2$) & $0.0137$ & Lateral Malleolus $\to$ Distal Phalanges \\
\hline
\end{tabular}
\end{table}

\subsection{Global Center of Gravity ($\mathbf{P}_{\text{CoG}}$)}
For each video frame $t$, the global 2D Center of Gravity $(X_{\text{CoG}}, Y_{\text{CoG}})$ is computed as the mass-weighted sum of individual segment centers of mass $\mathbf{p}_i = (x_i, y_i)$:
\begin{equation}
    X_{\text{CoG}}(t) = \sum_{i=1}^{S} w_i \cdot x_i(t), \quad Y_{\text{CoG}}(t) = \sum_{i=1}^{S} w_i \cdot y_i(t)
    \label{eq:cog_formula}
\end{equation}
where $\sum_{i=1}^{S} w_i = 1.0$.

\subsection{Stance Depth ($\mathcal{D}_{\text{stance}}$) \& Base of Support}
Stance depth is extracted as the Euclidean distance between the lead and rear ankle joint coordinates projected onto the horizontal ground plane:
\begin{equation}
    \mathcal{D}_{\text{stance}}(t) = \left| x_{\text{ankle, lead}}(t) - x_{\text{ankle, rear}}(t) \right| \cdot S_{px}^{-1}
\end{equation}
where $S_{px}$ is the pixel-to-meter scaling factor obtained from the calibration phase.
